\section{Conclusion}

\label{sec:conclusion}

We have presented a threshold variant of TFHE that enables boolean circuit evaluation
on encrypted data with distributed decryption. The CMPCOMBINE optimization reduces
multi-bit comparison from $2k-1$ to $k$ bootstrappings, making encrypted integer
operations practical. The threshold decryption protocol preserves IND-CPA security
under the LWE assumption while requiring only one round of communication. Integration
with the Lux FHE VM enables on-chain confidential computation with the full
expressiveness of boolean circuits.

\begin{thebibliography}{99}
\bibitem{chillotti2020tfhe} I. Chillotti, N. Gama, M. Georgieva, and M. Izabach\`ene, ``TFHE: Fast fully homomorphic encryption over the torus,'' \emph{J. Cryptology}, 2020.
\bibitem{gentry2009fully} C. Gentry, ``Fully homomorphic encryption using ideal lattices,'' in \emph{STOC}, 2009.
\bibitem{regev2009lattices} O. Regev, ``On lattices, learning with errors, random linear codes, and cryptography,'' \emph{JACM}, 2009.
\bibitem{brakerski2014leveled} Z. Brakerski, C. Gentry, and V. Vaikuntanathan, ``(Leveled) fully homomorphic encryption without bootstrapping,'' \emph{ACM TOCT}, 2014.
\bibitem{boneh2018threshold} D. Boneh et al., ``Threshold cryptosystems from threshold fully homomorphic encryption,'' in \emph{CRYPTO}, 2018.
\bibitem{asharov2012multiparty} G. Asharov et al., ``Multiparty computation with low communication, computation and interaction via threshold FHE,'' in \emph{EUROCRYPT}, 2012.
\end{thebibliography}

